\documentclass[aip,jcp,reprint,noshowkeys,superscriptaddress]{revtex4-1}
\usepackage{graphicx,dcolumn,bm,xcolor,microtype,multirow,amscd,amsmath,amssymb,amsfonts,physics,longtable,wrapfig}
\usepackage[version=4]{mhchem}

\usepackage[utf8]{inputenc}
\usepackage[T1]{fontenc}
\usepackage{txfonts}

\usepackage[
    colorlinks,
    linkcolor={red!50!black},
    citecolor={red!70!black},
    urlcolor={red!80!black}
	]{hyperref}
\urlstyle{same}

\newcommand{\ie}{\textit{i.e.}}
\newcommand{\eg}{\textit{e.g.}}
\newcommand{\alert}[1]{\textcolor{red}{#1}}
\usepackage[normalem]{ulem}
\newcommand{\titou}[1]{\textcolor{red}{#1}}
\newcommand{\trashPFL}[1]{\textcolor{\red}{\sout{#1}}}
\newcommand{\PFL}[1]{\titou{(\underline{\bf PFL}: #1)}}

\newcommand{\cre}[1]{\Hat{a}_{#1}^{\dag}}
\newcommand{\ani}[1]{\Hat{a}_{#1}}
\newcommand{\ca}[2]{\Hat{a}^{#1}_{#2}}

\newcommand{\mc}{\multicolumn}
\newcommand{\fnm}{\footnotemark}
\newcommand{\fnt}{\footnotetext}
\newcommand{\tabc}[1]{\multicolumn{1}{c}{#1}}
\newcommand{\QP}{\textsc{quantum package}}
\newcommand{\T}[1]{#1^{\intercal}}

% coordinates
\newcommand{\br}{\boldsymbol{r}}
\newcommand{\bx}{\boldsymbol{x}}
\newcommand{\bI}{\boldsymbol{1}}
\newcommand{\cK}{\mathcal{K}}

% orbital energies
\newcommand{\eps}{\epsilon}
\newcommand{\irr}{\text{irr}}

% shortcuts for greek letters
\newcommand{\si}{\sigma}
\newcommand{\la}{\lambda}
\newcommand{\ii}{\mathrm{i}}
\newcommand{\up}{\uparrow}
\newcommand{\dw}{\downarrow}
\newcommand{\upup}{\up\up}
\newcommand{\updw}{\up\dw}
\newcommand{\dwup}{\dw\up}
\newcommand{\dwdw}{\dw\downarrow}

\newcommand{\h}{\text{1h}}
\newcommand{\p}{\text{1p}}
\newcommand{\hp}{\text{1h+1p}}
\newcommand{\hhp}{\text{2h1p}}
\newcommand{\pph}{\text{2p1h}}
\newcommand{\ph}{\text{ph}}
\newcommand{\eh}{\text{eh}}
\newcommand{\he}{\text{he}}
\newcommand{\ee}{\text{ee}}
\newcommand{\teh}{\overline{\text{eh}}}
\newcommand{\pp}{\text{pp}}
\newcommand{\hh}{\text{hh}}
\newcommand{\xc}{\text{xc}}
\newcommand{\Hx}{\text{Hx}}

% addresses
\newcommand{\LCPQ}{Laboratoire de Chimie et Physique Quantiques (UMR 5626), Universit\'e de Toulouse, CNRS, UPS, France}

\begin{document}

\title{Third-Order Extension of the Multiresolvent Representation}

\author{Pierre-Fran\c{c}ois \surname{Loos}}
    \email{loos@irsamc.ups-tlse.fr}
    \affiliation{\LCPQ}

\begin{abstract}
These companion notes extend the multiresolvent representation of the parquet
vertex to third order in the bare interaction.
The frequency conventions, channel decomposition, first frequency reduction,
and complete second-order construction are taken from the second-order notes
and are not repeated here.
At third order, the standard $\eh$ and $\pp$ central spaces generate three
contributions: a second-order correction to the left coupling vertex, a
second-order correction to the right coupling vertex, and a first-order
correction to the central resolvent.
An unrestricted evaluation of the frequency convolutions shows, however, that
a frequency-dependent second-order irreducible vertex also activates internal
orbital sectors that vanish for an instantaneous kernel: same-occupancy
virtual--virtual and occupied--occupied sectors in the $\eh$ channel and
mixed occupied--virtual sectors in the $\pp$ channel.
We keep these additional terms explicitly and identify where the auxiliary
spaces of the second-order representation must be enlarged.
The purpose of these notes is therefore to isolate the genuinely new
third-order structures while keeping the notation and organization of the
second-order construction.
\end{abstract}

\maketitle

%%%%%%%%%%%%%%%%%%%%%%%%%
\section{Scope and conventions}
\label{sec:introduction}
%%%%%%%%%%%%%%%%%%%%%%%%%

These notes begin where the second-order multiresolvent construction ends.
We adopt without rederivation the direct-$\eh$ frequency convention, parquet
channel definitions, normalized factors $Q^\pm$, channel metrics and
resolvents, coupling vertices $U^r$ and $V^r$, and the explicit second-order
irreducible kernels $\Gamma^{\eh,(2)}$ and $\Gamma^{\pp,(2)}$.
The second-order derivation and its independent consistency checks are not
repeated.
The objective here is solely to identify and organize the structures that
first appear at third order.

As in the second-order notes, we use a Hartree--Fock zeroth-order reference
with a M{\o}ller--Plesset partition.
Accordingly, $G_0$ and $L_0$ denote the corresponding non-interacting
one- and two-particle propagators.
All expressions are written for real-valued orbitals, and causal
prescriptions are denoted by $0^+$.

\section{Third-order multiresolvent extension}
\label{sec:multiresolvent_third}
%%%%%%%%%%%%%%%%%%%%%%%%%

We take the complete second-order multiresolvent construction as given and
extend it by one perturbative order.
Only the generic channel factorization is needed below,
\begin{equation}
    \ii\Phi^r
    =
    U^r R^r V^r,
    \qquad
    R^r(\Omega_r)
    =
    [\Omega_r\eta^r-M^r]^{-1},
    \qquad
    r=\eh,\pp
    \label{eq:third_starting_point}
\end{equation}
with the frequency conventions, metrics, normalized one-frequency factors,
and all zeroth- and first-order quantities defined in the second-order notes.
Within the conventional central channel spaces, the third-order terms are
organized by corrections to the two coupling vertices and to the central
resolvent.
The unrestricted frequency convolution will show that additional orbital
sectors are activated at this order; these are introduced explicitly below
rather than being hidden inside the restricted ansatz.

For either primitive channel $r=\eh,\pp$, we expand
\begin{align}
	U^r
	&=
	U^{r,(1)}+U^{r,(2)}+\mathcal O(v^3)
	\\
	V^r
	&=
	V^{r,(1)}+V^{r,(2)}+\mathcal O(v^3)
	\label{eq:target_UV_expansion}
	\\
	M^r
	&=
	M^{r,(0)}+M^{r,(1)}+\mathcal O(v^2)
	\label{eq:target_M_expansion}
\end{align}
The corresponding central resolvent becomes
\begin{align}
	R^r
	&=
	R^{r,(0)}
	+
	R^{r,(1)}
	+
	\mathcal O(v^2)
	\label{eq:target_R_expansion}
	\\
	R^{r,(1)}
	&=
	R^{r,(0)}M^{r,(1)}R^{r,(0)}
	\label{eq:target_R_first}
\end{align}
The third-order contribution generated within the standard central spaces is
\begin{equation}
	\begin{split}
	\ii\Phi_{\rm std}^{r,(3)}
	={}&
	U^{r,(2)}R^{r,(0)}V^{r,(1)}
	+
	U^{r,(1)}R^{r,(0)}V^{r,(2)}
	\\
	&+
	U^{r,(1)}R^{r,(1)}V^{r,(1)}
	\end{split}
	\label{eq:target_third_order}
\end{equation}
The three terms correspond, respectively, to a second-order correction to
the left coupling vertex, a second-order correction to the right coupling
vertex, and a first-order correction to the central resolvent.
They exhaust the contribution generated within the central spaces of the
second-order construction, but not the unrestricted third-order convolution.
The additional orbital-sector contributions are identified explicitly in the
following sections.

\section{Direct eh vertex through third order}
\label{sec:eh_third}
%%%%%%%%%%%%%%%%%%%%%%%%%

The direct-$\eh$ construction uses without repetition the second-order
irreducible kernel $\Gamma^{\eh,(2)}$, the normalized factors $Q^\pm$,
the zeroth-order neutral resolvent $R^{\eh,(0)}$, and the first-order
coupling vertices $U^{\eh,(1)}$ and $V^{\eh,(1)}$ established in the
second-order notes.
The new ingredients are the second-order coupling-vertex corrections, the
first-order correction to the neutral resolvent, and the orbital sectors
that appear only in the unrestricted third-order convolution.

%%%%%%%%%%%%%%%%%%%%%%%%%
\subsection{Perturbative BSE expansion}
%%%%%%%%%%%%%%%%%%%%%%%%%

At third order, the direct-$\eh$ reducible vertex contains three
topologically distinct contributions.
Suppressing the explicit internal-frequency integrations with the convolution
symbol $\circ$, and retaining the same phase convention as in the
second-order notes, they can be written as
\begin{equation}
	\begin{split}
	\ii\Phi^{\eh,(3)}
	&=
	\left[\ii\Gamma^{\eh,(2)}\right]
	\circ\left[-\ii L_0\right]
	\circ\left[\ii\Gamma^{\eh,(1)}\right]
	\\
	&+
	\left[\ii\Gamma^{\eh,(1)}\right]
	\circ\left[-\ii L_0\right]
	\circ\left[\ii\Gamma^{\eh,(2)}\right]
	\\
	&+
	\left[\ii\Gamma^{\eh,(1)}\right]
	\circ\left[-\ii L_0\right]
	\circ\left[\ii\Gamma^{\eh,(1)}\right]
	\circ\left[-\ii L_0\right]
	\circ\left[\ii\Gamma^{\eh,(1)}\right]
	\end{split}
	\label{eq:Phi_eh_third_convolution}
\end{equation}
The first two terms contain the frequency-dependent second-order irreducible
kernel on the left or right of the central propagation.
The last term is the three-interaction $\eh$ ladder.

For a two-body interaction, the fully two-particle-irreducible vertex has no
second- or third-order contribution beyond the bare interaction,
\begin{equation}
	\ii\Lambda^{(1)}_{pqrs}
	=
	\mel{pq}{}{rs}
	\qquad
	\Lambda
	=
	\Lambda^{(1)}+\mathcal O(v^4)
	\label{eq:Lambda_through_third}
\end{equation}
so that all third-order contributions considered here are generated from
the channel-reducible pieces.

\subsection{Orbital sectors in the third-order convolution}
%%%%%%%%%%%%%%%%%%%%%%%%%

The orbital restriction implicit in the ordinary $\eh$ central resolvent
requires special care at third order.
For an instantaneous interaction, same-occupancy non-interacting bubbles
vanish,
\begin{equation}
	\int\frac{\dd{x}}{2\pi}
	G_{0,a}(x+\omega)G_{0,b}(x)
	=
	0
	\qquad
	\int\frac{\dd{x}}{2\pi}
	G_{0,i}(x+\omega)G_{0,j}(x)
	=
	0
	\label{eq:same_occupancy_bubbles_zero}
\end{equation}
because the two one-particle poles lie in the same half-plane.
This is why the second-order direct-$\eh$ construction can be represented
entirely in the forward and backward electron-hole sectors.

A frequency-dependent second-order irreducible vertex changes this
conclusion.
For example, with $a,b,c$ virtual and $k$ occupied,
\begin{equation}
	\begin{split}
	&\int\frac{\dd{x}}{2\pi}
	\frac{1}{x+\omega-\eps_a+\ii0^+}
	\frac{1}{x-\eps_b+\ii0^+}
	\frac{1}{x-\nu+(\eps_c-\eps_k)-\ii0^+}
	\\
	&\qquad=
	\frac{\ii}
	{\nu+\omega-(\eps_a+\eps_c-\eps_k)+\ii0^+}
	\frac{1}
	{\nu-(\eps_b+\eps_c-\eps_k)+\ii0^+}
	\end{split}
	\label{eq:vv_dynamic_counterexample}
\end{equation}
which is nonzero.
The analogous occupied--occupied contributions are also nonzero when the
dynamic kernel supplies a pole in the opposite half-plane.

It is therefore useful to separate the non-interacting two-particle
propagator into the conventional electron-hole sectors and the
same-occupancy sectors,
\begin{equation}
	L_0
	=
	L_{0,\eh}
	+
	L_{0,\rm so}
	\label{eq:L0_sector_decomposition}
\end{equation}
where $L_{0,\eh}$ contains the $ia$ and $ai$ sectors and
$L_{0,\rm so}$ contains the $ab$ and $ij$ sectors.
The $Q^\pm$ projections of the second-order construction account exactly for
the $L_{0,\eh}$ contribution.
The remaining third-order terms are defined by
\begin{align}
	\Delta_L^{\eh,(3)}
	&=
	\left[\ii\Gamma^{\eh,(2)}\right]
	\circ\left[-\ii L_{0,\rm so}\right]
	\circ\left[\ii\Gamma^{\eh,(1)}\right]
	\label{eq:Delta_eh_left_same_occ}
	\\
	\Delta_R^{\eh,(3)}
	&=
	\left[\ii\Gamma^{\eh,(1)}\right]
	\circ\left[-\ii L_{0,\rm so}\right]
	\circ\left[\ii\Gamma^{\eh,(2)}\right]
	\label{eq:Delta_eh_right_same_occ}
\end{align}
The purely first-order ladder has no analogous contribution because the
interaction is frequency independent and the same-occupancy bubble vanishes
before the next scattering event.
Consequently, the ordinary electron-hole central resolvent gives the complete
third-order contribution within the $ia/ai$ space, while
Eqs.~\eqref{eq:Delta_eh_left_same_occ} and
\eqref{eq:Delta_eh_right_same_occ} must be retained for the unrestricted
three-frequency vertex.

\subsection{Second-order corrections to the coupling vertices}
%%%%%%%%%%%%%%%%%%%%%%%%%

The left coupling-vertex correction for forward ($+$) central electron-hole propagation through $ia$ is obtained by integrating over one
fermionic frequency of the second-order irreducible kernel,
\begin{equation}
	U^{\eh,+,(2)}_{pr,ia}(\nu,\omega)
	=
	\frac{1}{2\pi}
	\int \dd{\bar\nu}
	\left[
	\ii\Gamma^{\eh,(2)}_{pa\,ri}
	(\nu,\bar\nu,\omega)
	\right]
	Q^+_{ia}(\bar\nu,\omega)
	\label{eq:U_eh_second_def}
\end{equation}
The explicit result is
\begin{equation}
	\begin{split}
	U^{\eh,+,(2)}_{pr,ia}(\nu,\omega)
	&=
	\sum_{jb}
	\frac{
	\mel{pb}{}{ij}
	\mel{aj}{}{rb}
	}
	{\nu-(\eps_i+\eps_j-\eps_b)-\ii0^+}
	\\
	&-
	\sum_{jb}
	\frac{
	\mel{pj}{}{ib}
	\mel{ab}{}{rj}
	}
	{\nu+\omega-(\eps_a+\eps_b-\eps_j)+\ii0^+}
	\\
	&-
	\frac{1}{2}
	\sum_{jk}
	\frac{
	\mel{pa}{}{jk}
	\mel{jk}{}{ri}
	}
	{\nu-(\eps_j+\eps_k-\eps_a)-\ii0^+}
	\\
	&+
	\frac{1}{2}
	\sum_{bc}
	\frac{
	\mel{pa}{}{bc}
	\mel{bc}{}{ri}
	}
	{\nu+\omega-(\eps_b+\eps_c-\eps_i)+\ii0^+}
	\end{split}
	\label{eq:U_eh_second_explicit}
\end{equation}
The first two terms originate from the transverse $\eh$ channel
and the last two from the $\hh$ and $\pp$ branches.

The right coupling-vertex correction for forward ($+$) central electron-hole propagation through $ia$ follows from the right side of the $\eh$ BSE with its orbital ordering kept explicit,
\begin{equation}
	V^{\eh,+,(2)}_{ia,sq}(\nu',\omega)
	=
	\frac{1}{2\pi}
	\int \dd{\bar\nu}
	Q^+_{ia}(\bar\nu,\omega)
	\left[
	\ii\Gamma^{\eh,(2)}_{iq\,as}
	(\bar\nu,\nu',\omega)
	\right]
	\label{eq:V_eh_second_def}
\end{equation}
which gives
\begin{equation}
	\begin{split}
	V^{\eh,+,(2)}_{ia,sq}(\nu',\omega)
	&=
	\sum_{kc}
	\frac{
	\mel{sc}{}{ik}
	\mel{ak}{}{qc}
	}
	{\nu'-(\eps_i+\eps_k-\eps_c)-\ii0^+}
	\\
	&-
	\sum_{kc}
	\frac{
	\mel{sk}{}{ic}
	\mel{ac}{}{qk}
	}
	{\nu'+\omega-(\eps_a+\eps_c-\eps_k)+\ii0^+}
	\\
	&-
	\frac{1}{2}
	\sum_{kl}
	\frac{
	\mel{sa}{}{kl}
	\mel{kl}{}{qi}
	}
	{\nu'-(\eps_k+\eps_l-\eps_a)-\ii0^+}
	\\
	&+
	\frac{1}{2}
	\sum_{cd}
	\frac{
	\mel{sa}{}{cd}
	\mel{cd}{}{qi}
	}
	{\nu'+\omega-(\eps_c+\eps_d-\eps_i)+\ii0^+}
	\end{split}
	\label{eq:V_eh_second_explicit}
\end{equation}
For the instantaneous first-order interaction, $\Gamma_{iq\,as}$ can be
related to alternative index groupings by ordinary integral symmetries.
For the frequency-dependent second-order kernel this replacement is not
permitted unless the corresponding frequency transformation is applied as
well; Eq.~\eqref{eq:V_eh_second_def} therefore fixes the right-side routing
used below.
The backward ($-$) coupling-vertex components follow from the same reversal operation as
the corresponding second-order projections.
The explicit formulas are not required for the third-order matching.

Equations~\eqref{eq:U_eh_second_explicit} and
\eqref{eq:V_eh_second_explicit} show that the second-order coupling
vertices retain exactly the $\hhp$ and $\pph$ pole spaces identified in
the second-order notes, now with the orbital grouping appropriate to the
left and right sides of the direct-$\eh$ reducible vertex.

%%%%%%%%%%%%%%%%%%%%%%%%%
\subsection{First-order correction to the neutral resolvent}
%%%%%%%%%%%%%%%%%%%%%%%%%

The first-order neutral interaction matrix is
\begin{equation}
	M^{\eh,(1)}
	=
	\begin{pmatrix}
	A^{(1)}&B^{(1)}
	\\
	B^{(1)\intercal}&A^{(1)\intercal}
	\end{pmatrix}
	\label{eq:M_eh_first}
\end{equation}
with
\begin{align}
	A^{(1)}_{ia,jb}
	&=
	\mel{aj}{}{ib}
	\label{eq:A_eh_first}
	\\
	B^{(1)}_{ia,jb}
	&=
	\mel{ab}{}{ij}
	\label{eq:B_eh_first}
\end{align}
Using the general expansion of Eq.~\eqref{eq:target_R_first}, the first-order neutral-resolvent correction is
\begin{equation}
	R^{\eh,(1)}
	=
	R^{\eh,(0)}
	M^{\eh,(1)}
	R^{\eh,(0)}
	\label{eq:R_eh_first}
\end{equation}
No additional metric matrix appears in Eq.~\eqref{eq:R_eh_first} because the
metric is already included in the definition of the resolvent.

Keeping only the forward electron-hole block for illustration, the
corresponding central third-order contribution is
\begin{equation}
	\begin{split}
	\left[
	\ii\Phi_{\rm cent}^{\eh,(3)}
	\right]_{pqrs}
	={}&
	\sum_{ia,jb}
	\frac{
	\mel{pa}{}{ri}
	\mel{aj}{}{ib}
	\mel{qj}{}{sb}
	}
	{
	[\omega-(\eps_a-\eps_i)+\ii0^+]
	[\omega-(\eps_b-\eps_j)+\ii0^+]
	}
	\end{split}
	\label{eq:Phi_eh_third_central_tda}
\end{equation}
The full forward-backward expression is generated automatically by the
first-order correction to the neutral resolvent
\begin{equation}
	U^{\eh,(1)}
	R^{\eh,(1)}
	V^{\eh,(1)}
	=
	U^{\eh,(1)}
	R^{\eh,(0)}
	M^{\eh,(1)}
	R^{\eh,(0)}
	V^{\eh,(1)}
	\label{eq:Phi_eh_third_central_matrix}
\end{equation}

%%%%%%%%%%%%%%%%%%%%%%%%%
\subsection{Direct eh result through third order}
%%%%%%%%%%%%%%%%%%%%%%%%%

Within the conventional forward/backward electron-hole space, the third-order
contribution has the same multiresolvent organization as at second order,
\begin{equation}
	\begin{split}
	\ii\Phi_{\rm std}^{\eh,(3)}
	={}&
	U^{\eh,(2)}R^{\eh,(0)}V^{\eh,(1)}
	+
	U^{\eh,(1)}R^{\eh,(0)}V^{\eh,(2)}
	\\
	&+
	U^{\eh,(1)}R^{\eh,(1)}V^{\eh,(1)}
	\end{split}
	\label{eq:Phi_eh_third_standard}
\end{equation}
The three terms correct, respectively, the left coupling vertex, the right
coupling vertex, and the central electron-hole resolvent.
The unrestricted result additionally contains the same-occupancy
contributions,
\begin{equation}
	\ii\Phi^{\eh,(3)}
	=
	\ii\Phi_{\rm std}^{\eh,(3)}
	+
	\Delta_L^{\eh,(3)}
	+
	\Delta_R^{\eh,(3)}
	\label{eq:Phi_eh_third_final}
\end{equation}

The multiresolvent character is already explicit in the first two terms.
For example, the forward left correction contains products of the form
\begin{equation}
	\frac{1}
	{\nu-(\eps_i+\eps_j-\eps_b)-\ii0^+}
	\frac{1}
	{\omega-(\eps_a-\eps_i)+\ii0^+}
	\label{eq:multi_resolvent_left_hhp}
\end{equation}
and
\begin{equation}
	\frac{1}
	{\nu+\omega-(\eps_a+\eps_b-\eps_j)+\ii0^+}
	\frac{1}
	{\omega-(\eps_a-\eps_i)+\ii0^+}
	\label{eq:multi_resolvent_left_pph}
\end{equation}
which combine an $(N-1)$ or $(N+1)$ coupling-vertex pole with the neutral
$N$-electron central resolvent.
The third term of Eq.~\eqref{eq:Phi_eh_third_standard} instead contains two
neutral resolvents separated by the first-order interaction matrix.

\section{pp vertex through third order}
\label{sec:pp_third}
%%%%%%%%%%%%%%%%%%%%%%%%%

The $\pp$ construction is organized in direct analogy with the $\eh$
channel.
We take the zeroth-order metric, central resolvent, and first-order coupling
vertices from the second-order notes and introduce only the ingredients that
first enter at third order.

The first-order correction to the central $\pp$ matrix is
\begin{equation}
	M^{\pp,(1)}
	=
	\begin{pmatrix}
	\mel{ab}{}{cd}&\mel{ab}{}{ij}
	\\
	\mel{ij}{}{cd}&\mel{ij}{}{kl}
	\end{pmatrix}
	\label{eq:M_pp_first}
\end{equation}
where the upper-left block acts in the $N+2$ sector, the lower-right block
in the $N-2$ sector, and the off-diagonal blocks couple the two sectors.
Accordingly,
\begin{equation}
	R^{\pp,(1)}
	=
	R^{\pp,(0)}
	M^{\pp,(1)}
	R^{\pp,(0)}
	\label{eq:R_pp_first}
\end{equation}

The second-order corrections $U^{\pp,(2)}$ and $V^{\pp,(2)}$ are obtained by
the same first frequency reduction as in the $\eh$ channel, now starting
from $\Gamma^{\pp,(2)}=\Phi^{\eh,(2)}+\Phi^{\teh,(2)}$.
Their poles again belong to the $(N-1)$ $\hhp$ and $(N+1)$ $\pph$ sectors,
with orbital-energy combinations of the form
$\eps_i+\eps_j-\eps_a$ and $\eps_a+\eps_b-\eps_i$.
The precise surviving fermionic frequency depends on which member of the
antisymmetrized $\pp$ external state is attached to the second-order kernel.
Since no new pole space is introduced, we do not introduce any additional
intermediate-state shorthand here.

Within the conventional $\ee/\hh$ central space, the third-order contribution
is therefore
\begin{equation}
	\begin{split}
	\ii\Phi_{\rm std}^{\pp,(3)}
	={}&
	U^{\pp,(2)}R^{\pp,(0)}V^{\pp,(1)}
	+
	U^{\pp,(1)}R^{\pp,(0)}V^{\pp,(2)}
	\\
	&+
	U^{\pp,(1)}R^{\pp,(1)}V^{\pp,(1)}
	\end{split}
	\label{eq:Phi_pp_third_standard}
\end{equation}
The central term contains the third-order $\pp$ and $\hh$ ladders together
with the corresponding forward/backward couplings.
For example, the pure $\ee$ contribution is
\begin{equation}
	\begin{split}
	\left[
	\ii\Phi_{\rm cent}^{\pp,(3)}
	\right]_{pqrs}^{\ee}
	={}&
	\sum_{\substack{a<b\\c<d}}
	\frac{
	\mel{pq}{}{ab}
	\mel{ab}{}{cd}
	\mel{cd}{}{rs}
	}
	{\Omega_\pp-(\eps_a+\eps_b)+\ii0^+}
	\\
	&\times
	\frac{1}
	{\Omega_\pp-(\eps_c+\eps_d)+\ii0^+}
	\end{split}
	\label{eq:Phi_pp_third_central_ee}
\end{equation}

As in the $\eh$ channel, the conventional central space is not sufficient
for the unrestricted convolution when the second-order irreducible vertex is
frequency dependent.
Mixed occupied--virtual internal sectors, which vanish for an isolated
instantaneous $\pp$ contraction, can then contribute.
We denote the corresponding left and right contributions by
$\Delta_L^{\pp,(3)}$ and $\Delta_R^{\pp,(3)}$.
Their explicit phase, normalization, and frequency routing depend on the
native antisymmetrized $\pp$ BSE convention and are not evaluated here.
The unrestricted result is therefore organized as
\begin{equation}
	\ii\Phi^{\pp,(3)}
	=
	\ii\Phi_{\rm std}^{\pp,(3)}
	+
	\Delta_L^{\pp,(3)}
	+
	\Delta_R^{\pp,(3)}
	\label{eq:Phi_pp_third_final}
\end{equation}
At this stage the $\pp$ extension should therefore be regarded as a formal
organization of the third-order contributions rather than a complete
implementable orbital expression.

In the direct-$\eh$ frequency convention,
$\Omega_\pp=\omega+\nu+\nu'$, so the central $\pp$ resolvent carries the
characteristic anti-diagonal bosonic frequency of this channel.

\section{Consistency checks at third order}
\label{sec:checks_third}
%%%%%%%%%%%%%%%%%%%%%%%%%

The third-order extension is deliberately overdetermined.
Several independent checks can therefore be applied to the new terms; the
second-order validations are not repeated here.

%%%%%%%%%%%%%%%%%%%%%%%%%
\subsection{Status of the second frequency reduction}
%%%%%%%%%%%%%%%%%%%%%%%%%

The second frequency reduction provides a useful check only after the
four-point object has been mapped unambiguously onto the corresponding
forward/backward channel blocks.
At third order this mapping must be applied separately to each left, right,
and central contribution, including the complementary orbital sectors.
A direct reuse of a second-order block formula without the associated index
and frequency transformation can interchange a forward denominator
$\omega-(\eps_a-\eps_i)$ with its reversed counterpart
$\omega+(\eps_a-\eps_i)$.

For this reason, we do not use a single projected one-frequency identity as
an independent validation of the third-order construction.
The algebraic resolvent relation
\begin{equation}
	R^{r,(1)}
	=
	R^{r,(0)}M^{r,(1)}R^{r,(0)}
	\qquad
	r=\eh,\pp
	\label{eq:R_first_repeated}
\end{equation}
is exact for the stated resolvent definition, while a complete second
frequency reduction requires the same explicit orbital and frequency mapping
for every term, including $\Delta_{L/R}^{r,(3)}$.

\subsection{Crossing}
%%%%%%%%%%%%%%%%%%%%%%%%%

The transverse $\eh$ primitive is never constructed
independently.
At every perturbative order,
\begin{equation}
	\Phi^{\teh,(n)}_{pqrs}(\nu,\nu',\omega)
	=
	-
	\Phi^{\eh,(n)}_{pqsr}
	\left(
	\nu,
	\nu+\omega,
	\nu'-\nu
	\right)
	\label{eq:crossing_order_n}
\end{equation}
This maps each left, right, and central third-order direct $\eh$
contribution into the corresponding transverse topology and preserves the
fermionic crossing symmetry order by order.

%%%%%%%%%%%%%%%%%%%%%%%%%
\subsection{Parquet reconstruction}
%%%%%%%%%%%%%%%%%%%%%%%%%

Because the fully irreducible vertex has no third-order correction,
\begin{align}
	F^{(3)}
	&=
	\Phi^{\eh,(3)}+\Phi^{\teh,(3)}+\Phi^{\pp,(3)}
	\label{eq:F_n_parquet}
	\\
	\Gamma^{\eh,(3)}
	&=
	\Phi^{\teh,(3)}+\Phi^{\pp,(3)}
	\label{eq:Gamma_eh_n}
	\\
	\Gamma^{\pp,(3)}
	&=
	\Phi^{\eh,(3)}+\Phi^{\teh,(3)}
	\label{eq:Gamma_pp_n}
	\\
	F^{(3)}
	&=
	\Gamma^{\eh,(3)}+\Phi^{\eh,(3)}
	=
	\Gamma^{\pp,(3)}+\Phi^{\pp,(3)}
	\label{eq:channel_consistency_n}
\end{align}
The direct $\eh$ Hedin kernel is therefore reconstructed from the
two primitive multiresolvent objects through
$\Gamma^\eh=\Lambda+\Phi^{\teh}+\Phi^\pp$.

%%%%%%%%%%%%%%%%%%%%%%%%%
\subsection{High-frequency structure}
%%%%%%%%%%%%%%%%%%%%%%%%%

The second-order corrections to the coupling vertices decay as the inverse of their fermionic frequency,
\begin{equation}
	U^{(2)}
	=
	\mathcal O(\nu^{-1})
	\qquad
	V^{(2)}
	=
	\mathcal O({\nu'}^{-1})
	\label{eq:coupling_high_frequency}
\end{equation}
When both fermionic frequencies are large, only the central
$\omega$-dependent contribution survives at third order.
The resulting hierarchy is therefore consistent with the standard
parquet asymptotic decomposition into a purely bosonic part and left and
right one-fermionic-frequency corrections.

The same argument applies in the $\pp$ channel with the central bosonic
frequency $\Omega_\pp$ held fixed.

%%%%%%%%%%%%%%%%%%%%%%%%%
\subsection{Diagram classes and orbital-space completeness}
%%%%%%%%%%%%%%%%%%%%%%%%%

Every third-order diagram reducible in the direct $\eh$ channel belongs to
one of the three classes displayed in
Eq.~\eqref{eq:Phi_eh_third_convolution}: a second-order irreducible vertex
on the left, a second-order irreducible vertex on the right, or three
first-order interactions separated by two non-interacting propagations.
The corresponding statement holds in the $\pp$ channel.

This topological classification prevents double counting between parquet
channels, but it does not justify restricting the internal orbital sums of
the first two classes to the standard central spaces.
Because $\Gamma^{(2)}$ is frequency dependent, same-occupancy sectors in
the $\eh$ channel and mixed occupied--virtual sectors in the $\pp$ channel
can survive the frequency integration.
The standard $URV$ expressions are therefore complete within the
second-order central spaces, while the unrestricted three-frequency vertex
also contains $\Delta_{L/R}^{r,(3)}$.
Crossing and parquet reconstruction remain valid once these additional
terms are included.

\section{Recipe through third order}
\label{sec:recipe}
%%%%%%%%%%%%%%%%%%%%%%%%%

Starting from the second-order quantities $U^{r,(1)}$, $V^{r,(1)}$,
$R^{r,(0)}$, and $\Gamma^{r,(2)}$, with $r=\eh,\pp$, the third-order
construction proceeds as follows.

First, obtain the second-order corrections $U^{r,(2)}$ and $V^{r,(2)}$
from the corresponding second-order channel-irreducible vertex using the
same first frequency reduction as at second order.
Their nontrivial poles remain in the $(N-1)$ $\hhp$ and $(N+1)$ $\pph$
sectors.

Second, construct the first-order correction $M^{r,(1)}$ to the central
matrix and
\begin{equation}
	R^{r,(1)}
	=
	R^{r,(0)}M^{r,(1)}R^{r,(0)}
	\label{eq:recipe_R_first}
\end{equation}
The central resolvent describes neutral $N$-electron propagation for
$r=\eh$ and $(N\pm2)$-electron propagation for $r=\pp$.

Third, inspect the internal orbital sectors before imposing the standard
central-space restriction.
The frequency-dependent second-order kernel activates same-occupancy
virtual--virtual and occupied--occupied sectors in the $\eh$ channel and
mixed occupied--virtual sectors in the $\pp$ channel.
Their contributions are denoted by
$\Delta_L^{r,(3)}$ and $\Delta_R^{r,(3)}$.

The complete third-order organization is then
\begin{equation}
	\begin{split}
	\ii\Phi^{r,(3)}
	={}&
	U^{r,(2)}R^{r,(0)}V^{r,(1)}
	+
	U^{r,(1)}R^{r,(0)}V^{r,(2)}
	\\
	&+
	U^{r,(1)}R^{r,(1)}V^{r,(1)}
	+
	\Delta_L^{r,(3)}
	+
	\Delta_R^{r,(3)}
	\end{split}
	\label{eq:recipe_third_general}
\end{equation}
Generate $\Phi^{\teh,(3)}$ from $\Phi^{\eh,(3)}$ by crossing and reconstruct
\begin{align}
	\Gamma^{\eh,(3)}
	&=
	\Phi^{\teh,(3)}+\Phi^{\pp,(3)}
	\\
	F^{(3)}
	&=
	\Phi^{\eh,(3)}+\Phi^{\teh,(3)}+\Phi^{\pp,(3)}
\end{align}

The first three terms in Eq.~\eqref{eq:recipe_third_general} are the direct
third-order continuation of the second-order multiresolvent representation.
The $\Delta_{L/R}^{r,(3)}$ terms identify precisely where the standard
second-order central spaces cease to span the unrestricted third-order
frequency convolution.


\end{document}
